%% ================================================================================================
%%  5. Aggregating full ballots on an open agenda
%%  Notes: research-agenda.md §5; code: simulations/python/voting.py (rules), examples.py (E2)
%% ================================================================================================
\section{Aggregating full ballots on an open agenda}\label{sec:aggregation}

Under H2 every voter ranks every option, so the procedure needs a rule that aggregates full rankings. This
section compares the main families of rules on the properties that matter under H1--H3.

\subsection{The families of rules}

\begin{itemize}
  \item \textbf{Positional rules} give points by rank: plurality (one point to the first choice), Borda
  ($m-k$ points to the $k$-th choice; \citealp{Young1974}).
  \item \textbf{Runoff rules}: instant-runoff voting (IRV) eliminates the option with the fewest first
  choices until one has a majority.
  \item \textbf{Condorcet methods} elect the Condorcet winner when there is one and differ in how they
  resolve cycles. Copeland counts pairwise wins; Minimax (Simpson--Kramer) minimises the worst pairwise
  defeat \citep{Simpson1969}; Kemeny--Young finds the ranking that disagrees least with the voters'
  pairwise preferences \citep{Kemeny1959,YoungLevenglick1978}; Ranked Pairs locks pairwise majorities
  from the strongest down, skipping any that would close a cycle \citep{Tideman1987}; the Schulze method
  elects the option whose strongest chains of pairwise victories against every other option are at
  least as strong as the reverse chains \citep{Schulze2011}; Split Cycle discards the weakest majority in
  every cycle and elects the options left undefeated \citep{HollidayPacuit2023splitcycle}, and Stable
  Voting refines it to a single winner \citep{HollidayPacuit2023stable}.
  \item \textbf{Evaluative rules} ask for grades rather than rankings: approval voting elects the option
  the most voters find acceptable \citep{BramsFishburn1978,LaslierSanver2010}; Majority Judgment elects
  the option with the best median grade \citep{BalinskiLaraki2007,BalinskiLaraki2011}; see
  \citet{FelsenthalMachover2008} for a critical evaluation.
  \item \textbf{Probabilistic rules}: maximal lotteries choose a lottery over options that no other
  lottery beats in expected pairwise majority \citep{Fishburn1984,LaffondLaslierLeBreton1993}; they are
  the only rule with a set of strong consistency properties \citep{BrandlBrandtSeedig2016}, but a group
  choosing a single option by lot is a hard sell for most organisations.
  \item \textbf{Compromise rules} with vetoes, reviewed in Section~\ref{sec:veto}: fallback bargaining,
  the majoritarian compromise, voting by veto and the proportional veto core.
\end{itemize}
Kemeny--Young and Dodgson's rule are NP-hard to compute \citep{BartholdiToveyTrick1989hard}; all the
other rules above run in polynomial time, Schulze in $O(m^3)$ after the $O(nm^2)$ count of pairwise
preferences. \citet{ElkindFaliszewskiSlinko2015} show that many of these rules can be read as finding the
nearest ``consensus'' profile (unanimous, or with a Condorcet winner) under some distance between
profiles, which makes the choice of a rule a choice of what counts as consensus. \citet{BrandtEtAl2016handbook}
survey the axiomatic and computational results on all these rules.

\begin{table}[htbp]\centering\small
\caption{Properties of the rules compared}\label{tab:properties}
\footnotesize
\begin{tabular}{l c c c c c c}
\toprule
Rule & \shortstack{Condorcet\\winner} & \shortstack{Independent\\of clones} & Monotonic & Polynomial & \shortstack{Uses\\vetoes} & \shortstack{Never\\stuck} \\
\midrule
Plurality              & $\times$     & $\times$     & \checkmark & \checkmark & $\times$   & \checkmark \\
IRV                    & $\times$     & \checkmark   & $\times$   & \checkmark & $\times$   & \checkmark \\
Borda                  & $\times$     & $\times$     & \checkmark & \checkmark & $\times$   & \checkmark \\
Copeland               & \checkmark   & $\times$     & \checkmark & \checkmark & $\times$   & \checkmark \\
Minimax                & \checkmark   & $\times$     & \checkmark & \checkmark & $\times$   & \checkmark \\
Kemeny--Young          & \checkmark   & $\times$     & \checkmark & $\times$   & $\times$   & \checkmark \\
Ranked Pairs           & \checkmark   & \checkmark   & \checkmark & \checkmark & $\times$   & \checkmark \\
Schulze                & \checkmark   & \checkmark   & \checkmark & \checkmark & $\times$   & \checkmark \\
Split Cycle$^a$        & \checkmark   & \checkmark   & \checkmark & \checkmark & $\times$   & \checkmark \\
\midrule
Approval               & $\times$     & \checkmark$^b$ & \checkmark & \checkmark & \checkmark$^c$ & \checkmark \\
Majority Judgment      & $\times$     & \checkmark$^b$ & \checkmark & \checkmark & \checkmark$^c$ & \checkmark \\
\midrule
Fallback bargaining    & $\times$     & $\times$     & \checkmark & \checkmark & $\times$   & \checkmark \\
\quad with impasse     & $\times$     & $\times$     & \checkmark & \checkmark & \checkmark & $\times$   \\
Majoritarian compromise & $\times$    & $\times$     & \checkmark & \checkmark & $\times$   & \checkmark \\
Unanimity, status quo default & $\times$ & \checkmark & \checkmark & \checkmark & \checkmark & $\times$ \\
Schulze within the PVC & $\times$ & $\times$ & --- & \checkmark & \checkmark$^d$ & \checkmark \\
\QVC($\theta$), Section~\ref{sec:procedure} & $\times^e$ & \checkmark & \checkmark$^f$ & \checkmark & \checkmark & \checkmark \\
\bottomrule
\end{tabular}
\figurenotes{Condorcet winner: elects it whenever it exists. Independent of clones: Definition~\ref{def:clones}.
Monotonic: raising the winner in some rankings never makes it lose. Never stuck: always selects an option
in a fixed number of rounds without a privileged default (Definition~\ref{def:decisive}). $\times$: fails
in some profiles; ---: not established, to our knowledge. $^a$Not resolute: ties are frequent and need a
tie-break. $^b$When voters grade or approve clones alike. $^c$The lowest grade, or non-approval, acts as a
veto mark but blocks nothing. $^d$PVC: proportional veto core; veto power read from the bottom of the rankings. $^e$Elects the Condorcet
winner of the admissible options (Proposition~\ref{prop:condorcet-admissible}). $^f$For changes in the
rankings that leave the veto marks unchanged.}{\citet{Tideman1987,Schulze2011,HollidayPacuit2023splitcycle,
BramsKilgour2001,Felsenthal2012review,BartholdiToveyTrick1989hard}; Propositions~\ref{prop:pvc-clones}
and~\ref{prop:decisive}--\ref{prop:qvc-clones}; simulations of Table~\ref{tab:clones}.}
\end{table}

Table~\ref{tab:properties} summarises the properties; \citet{Felsenthal2012review} reviews the paradoxes
that afflict each rule in more detail. Three properties sort the candidates for our use.

\subsection{An open agenda calls for independence of clones}\label{sec:clones}

When anyone can propose options, the composition of the agenda becomes a strategy. Adding variants of an
option, its clones, affects rules in three ways \citep{Tideman1987}: under plurality the clones split
their supporters' votes and lose (vote splitting); under Borda's rule they push other options down
everyone's ranking and gain (teaming); under some other rules they can crowd out third options. A
faction can therefore flood the agenda with variants of its favourite, or, under plurality, sink an
option by proposing its near-copies. Ranked Pairs, the Schulze method and Split Cycle are independent of
clones \citep{Tideman1987,ZavistTideman1989,Schulze2011,HollidayPacuit2023splitcycle}; Copeland, Minimax
and Kemeny--Young are not.

Veto-based compromise rules are exposed too, because the number of options enters their definition. In
fallback bargaining, clones of the majority's favourite fill the top of the majority's rankings and delay
the depth at which the compromise appears on every list. In the proportional veto core, a coalition's
veto power is a share of the number of options, so each clone added by the majority costs the minority a
veto that it can no longer spend elsewhere.

\begin{proposition}[Flooding defeats proportional veto]\label{prop:pvc-clones}
Neither the proportional veto core nor fallback bargaining is independent of clones: there are profiles
in which adding clones of an option outside the core (respectively, not chosen) brings that option into
the core (respectively, gets it chosen).
\end{proposition}

Example~E2 of Section~\ref{sec:examples} gives such a profile: in Example~E1, the 49\% minority blocks the
majority's favourite $A$, which is outside the core; once the majority proposes two variants of $A$, the
minority's three vetoes no longer cover them all, $A$ enters the core, and Schulze within the core elects
$A$ (Table~\ref{tab:examples}). In our simulations, three clones of a losing option change the winner of
the proportional veto core with Schulze in 15\% of polarised elections and in 54\% under the impartial
culture; they change the outcome of fallback bargaining in 25--36\% of elections, of the majoritarian
compromise in 23--52\% and of Borda's rule in 14--24\%, but never that of Schulze, Ranked Pairs or Split
Cycle (Table~\ref{tab:clones}).

Independence of clones also answers the agenda-setting concern of Section~\ref{sec:veto}. With all
options on one ballot, nobody controls the order of votes; with a clone-independent rule, nobody gains
by inflating the agenda. H1 and H2 together then neutralise most of the power that proposal rights confer
in sequential procedures \citep{BaronFerejohn1989,BarberaGerber2017}.

\subsection{Majority efficiency, manipulation and welfare}

Condorcet methods are the natural candidates when the group wants the majority principle respected
whenever it can be. Among them, Ranked Pairs and Schulze are Condorcet- and Smith-consistent, independent
of clones, monotonic and computable for any realistic number of options. Schulze is used by many organisations, Debian among them
\citep{Debian2024constitution}; Split Cycle adds stronger protections against spoilers and no-show effects at the cost of
frequent ties \citep{HollidayPacuit2023splitcycle}.

On manipulation, \citet{GreenArmytageTidemanCosman2016} compare 54 rules on survey data, spatial models
and random cultures: IRV and its Condorcet hybrid (elect the Condorcet winner, else run IRV) are the
least manipulable, while most rules have very similar utilitarian efficiency, with plurality,
anti-plurality and random dictatorship clearly worse. Adding a Condorcet step to a rule almost never
creates new opportunities for manipulation. Plurality is the rule experts like least \citep{Laslier2012loser}.
Section~\ref{sec:simulation} repeats such a comparison for our setting, with vetoes and an open agenda.
